%% kapa: exam=1 round=28 year=2017 period=1 form=A subject=economics subject_ko=경제학원론 seq=1 number=41
%% source: Q-Net 감정평가사 기출문제 (cst00309) artlSeq=5199603 file=extracted/1008/5199603/2017년 제28회 감정평가사 1차시험 문제지/1교시 A형.pdf
%% conversion: kordoc 4.24.51 -> tools/gen1.py (auto); review=proofread-sonnet; flags=none
\begin{question}{41}
소득이 600인 소비자 甲은 X재와 Y재만을 소비하며 효용함수는 $U=x+y$이다. $P_{X}=20$, $P_{Y}=15$ 이던 두 재화의 가격이 $P_{X}=20$, $P_{Y}=25$로 변할 때 최적 소비에 관한 설명으로 옳은 것은? (단, $x$는 X재 소비량, $y$는 Y재 소비량이다.)
\begin{choices}
  \item X재 소비를 30 단위 증가시킨다.
  \item X재 소비를 40 단위 증가시킨다.
  \item Y재 소비를 30 단위 증가시킨다.
  \item Y재 소비를 40 단위 증가시킨다.
  \item Y재 소비를 30 단위 감소시킨다.
\end{choices}
\end{question}
\answer{1}
